\section{Conclusion}\label{sec:conclusion}

We introduced a remedy decomposition of the misclassified minority class that sorts
each error into \emph{threshold-recoverable}, \emph{data-reducible} (epistemic), or
\emph{irreducible} (aleatoric), and used it to explain why class-imbalance
correction so often reduces to a decision-threshold change. Across 79 real datasets
the minority deficit is overwhelmingly threshold-recoverable (mean $69\%$, $74\%$ at
imbalance ratio $>3$): the classifier already ranks these instances correctly and
only the default cutoff is wrong. This accounts mechanistically for a documented but
unexplained regularity, namely that SMOTE-family oversampling does not improve ranking
and a threshold
move reproduces its benefit without synthetic generation or probability
distortion, and
recasts the long-observed accuracy/recall tradeoff as an operating-point effect:
such oversampling is an implicit, inefficient enactment of a threshold shift, which on
controlled data we show it achieves by generating across the overlap boundary. Once
every strategy is given the same out-of-fold threshold tuning, the large
default-threshold gap between the families collapses to within $\pm1$\,pp across
three base learners (six evaluated; naive Bayes shows a small but significant
generative edge of $1.6$\,pp, and the single undertrained MLP, whose ranking
imbalance itself corrupts, marks the account's boundary). The
decomposition also says \emph{when} the threshold is not enough: the small
data-reducible fraction is where minority data genuinely helps, the irreducible
fraction where no remedy in our menu does, though a strong non-generative default is
near-optimal in the threshold-recoverable majority. The practical message is simple:
diagnose the deficit, and for most of it, move the decision rather than oversample.
